\documentclass[11pt]{article}
\usepackage[utf8]{inputenc}
\usepackage[T1]{fontenc}
\usepackage{amsmath,amssymb,amsthm,mathtools}
\usepackage[margin=1in]{geometry}
\usepackage{enumitem}
\usepackage{xcolor}

\newcommand{\PR}{\textcolor{green!55!black}{\textbf{[P]}}}
\newcommand{\NM}{\textcolor{blue!70!black}{\textbf{[N]}}}
\newcommand{\OP}{\textcolor{red!75!black}{\textbf{[O]}}}
\newcommand{\GAP}{\textcolor{red!75!black}{\textsc{[gap]}}}

\theoremstyle{plain}
\newtheorem{theorem}{Teorema}
\newtheorem{lemma}{Lema}
\newtheorem{proposition}{Proposici\'on}
\newtheorem{conjecture}{Conjetura}
\theoremstyle{definition}
\newtheorem{remark}{Observaci\'on}
\newtheorem{definition}{Definici\'on}

\newcommand{\R}{\mathbb{R}}
\newcommand{\Lam}{\Lambda}
\newcommand{\half}{\tfrac12}
\newcommand{\El}{E_\lambda}
\newcommand{\kin}{K}
\newcommand{\vv}{V}
\newcommand{\pole}{P}
\newcommand{\fp}{\varphi_{\mathrm{polo}}}
\DeclareMathOperator{\Real}{Re}
\DeclareMathOperator{\ran}{ran}

\title{\textbf{Ataque al \GAP\ (i): el d\'eficit de Hardy vive en el modo--polo.\\
Reducci\'on de Schur $+$ desigualdad de Hardy mejorada}}
\author{Programa \texttt{riemann-program} --- documento de trabajo}
\date{\today}

\begin{document}
\maketitle

\begin{abstract}
El \GAP\ (i) era: probar que el d\'eficit cin\'etico--pozo est\'a \emph{confinado} a la
direcci\'on--polo $\fp$, i.e.\ $\kin-\vv\succeq0$ sobre $\fp^\perp$. Mostramos
(\S\ref{sec:schur}) que, v\'ia complemento de Schur para la perturbaci\'on de rango $1$
$\pole=\kappa|\fp\rangle\langle\fp|$, la positividad $A_\lambda\succeq0$ \emph{equivale}
exactamente a dos condiciones: \textbf{(i)} $\kin-\vv\succeq0$ sobre $\fp^\perp$ y
\textbf{(ii)} una desigualdad escalar de Schur. La condici\'on (i) es una
\textbf{desigualdad de Hardy mejorada}: el extremal de Hardy (donde la cota cin\'etica
es ajustada $c=2$) es \emph{precisamente} el modo uniforme/constante $=\fp$; al
proyectarlo fuera, aparece un \emph{resto positivo} (Brezis--V\'azquez), y la
desigualdad cin\'etico--pozo se vuelve estricta. As\'i el \GAP\ (i) \textbf{no} es
RH--equivalente: es la versi\'on no--local de una mejora de Hardy, incondicional. El
residuo \GAP\ es la no--localidad del s\'imbolo arquimediano (Lévy/CND).
\end{abstract}

\tableofcontents

% =====================================================================
\section{La reducci\'on de Schur (exacta)}\label{sec:schur}
% =====================================================================

Sea $B:=\kin-\vv$ (cin\'etica menos pozo, sim\'etrica, del lado aritm\'etico) y
$\pole=\kappa_\lambda|\fp\rangle\langle\fp|$ el t\'ermino--polo (rango $1$, positivo,
incondicional). La forma de Weil es $A_\lambda=B+\pole$. En la descomposici\'on
$\El=\R\fp\oplus\fp^\perp$ escribimos en bloques
\[
  B=\begin{pmatrix} b & r^* \\ r & M\end{pmatrix},
  \qquad
  b=\langle\fp,B\fp\rangle\in\R,\quad r=\Pi_\perp B\fp,\quad M=\Pi_\perp B\,\Pi_\perp,
\]
donde $\Pi_\perp$ es la proyecci\'on a $\fp^\perp$. Entonces
$A_\lambda=\begin{pmatrix} b+\kappa_\lambda & r^*\\ r & M\end{pmatrix}$.

\begin{lemma}[Reducci\'on de Schur]\label{lem:schur}
$A_\lambda\succeq0$ \emph{si y s\'olo si}
\begin{enumerate}[label=\textbf{(\roman*)},leftmargin=2.2em]
\item $M\succeq0$\quad($\kin-\vv$ es semidefinida positiva sobre $\fp^\perp$);
\item $r\in\ran(M)$\ \ y\ \ $b+\kappa_\lambda\;\ge\;r^*M^{+}r$\quad(condici\'on escalar de
Schur, $M^+$ pseudoinversa).
\end{enumerate}
\end{lemma}
\begin{proof}
Criterio de Schur para matrices sim\'etricas en bloques con bloque $(2,2)=M$
semidefinido. \qed
\end{proof}

\noindent Esto \emph{formaliza} la separaci\'on que la v\'ia \#9 sugiri\'o: toda la
positividad de Weil se parte en \textbf{(i)} positividad sobre el complemento del polo,
y \textbf{(ii)} una \'unica desigualdad escalar a lo largo del polo. El \GAP\ (i) es
exactamente la condici\'on (i).

% =====================================================================
\section{El \GAP\ (i) es una desigualdad de Hardy mejorada}\label{sec:hardy}
% =====================================================================

La condici\'on (i), $M\succeq0$, es
\begin{equation}\label{eq:gapi}
  \kin(g)\;\ge\;\vv_\lambda(g)\qquad\text{para todo }g\in\El\cap\fp^\perp .
\end{equation}

\paragraph{Por qu\'e la versi\'on sin restricci\'on falla.}
La desigualdad \emph{sin} proyectar, $\kin(g)\ge\vv_\lambda(g)\ \forall g\in\El$, es la
balanza de Weil con constante ajustada $c=2$, que las t\'ecnicas elementales \emph{no}
alcanzan (dan $c<2$; el ajuste $c=2$ es RH; cf.\ \texttt{A\_intento\_demostracion}). El
d\'eficit $c<2$ es real.

\begin{proposition}[El d\'eficit es el extremal de Hardy $=\fp$]\label{prop:extremal}
La desigualdad de tipo Hardy subyacente,
$\kin(g)=\frac1{2\pi}\int|\widehat g|^2(\Omega(t)-\Omega(0))\,dt\ \gtrsim\
c\int|g'|^2$, tiene como funci\'on \emph{extremal} (donde la cota es ajustada) el modo
de \emph{frecuencia cero} / uniforme. En la base modal de $A_\lambda$, ese extremal es
el estado fundamental, que por la v\'ia \#9 es justamente $\fp$ (la direcci\'on blanda,
$=$ densidad global $=$ polo). \NM\ (v\'ia \#9: la curvatura del d\'eficit es rango $1$
sobre $\fp$).
\end{proposition}

\begin{remark}[Mejora de Hardy: el resto de Brezis--V\'azquez]
Es un fen\'omeno cl\'asico que las desigualdades de Hardy se vuelven \emph{estrictas con
resto positivo} al restringir al complemento ortogonal del extremal:
\[
  \int|g'|^2-c\!\int\!\frac{|g|^2}{\text{(peso)}}
  \;\ge\;\beta\int|g|^2_{\text{(resto)}}>0
  \qquad\text{para }g\perp(\text{extremal}).
\]
Si el extremal es $\fp$ (Prop.\ \ref{prop:extremal}), entonces sobre $\fp^\perp$ el
d\'eficit $c<2$ \emph{se cierra}: la cota se vuelve $c=2$ con margen, dando
\eqref{eq:gapi}. \textbf{El \GAP\ (i) es, por tanto, una mejora de Hardy, no RH.}
\end{remark}

\begin{conjecture}[Hardy mejorada no--local]\label{conj:hardy}
Para el operador cin\'etico no--local $\kin$ (s\'imbolo CND $\Omega$, medida de L\'evy
$d\nu(\xi)=2e^{-\xi/2}(1-e^{-2\xi})^{-1}d\xi$) y el pozo $\vv_\lambda$ (momentos PNT
$M_0,M_2$), existe $\beta_\lambda\ge0$ tal que
\[
  \kin(g)-\vv_\lambda(g)\;\ge\;\beta_\lambda\,\|g\|^2
  \qquad\text{para todo }g\in\El\cap\fp^\perp .
\]
\end{conjecture}

% =====================================================================
\section{Conteo de estados ligados: por qu\'e ``a lo sumo uno''}\label{sec:bargmann}
% =====================================================================

La Conjetura \ref{conj:hardy} equivale a que el ``Schr\"odinger'' no--local
$\kin-\vv_\lambda$ tenga \textbf{a lo sumo un estado ligado} (un autovalor negativo), y
que su autovector sea $\fp$.

\begin{remark}[CLR/Bargmann como herramienta incondicional]
El n\'umero de autovalores negativos de $\kin-\vv$ se acota por funcionales del pozo
(Cwikel--Lieb--Rozenblum, Bargmann, Lieb--Thirring) que dependen \emph{s\'olo} de
$\vv_\lambda$ (lado primo, PNT), \emph{no} de los ceros. Una cota
$N_-(\kin-\vv_\lambda)\le1$ ser\'ia incondicional. \textbf{Consistencia interna:}
$A_\lambda=(\kin-\vv)+\pole\succeq0$ con espectro $(k+1)^2\varepsilon_0>0$ (torre de
Weyl, \NM) implica que $\kin-\vv=A_\lambda-\pole$ tiene a lo sumo \emph{un} autovalor
negativo (resta de rango $1$), alineado con $\fp$ --- exactamente el cuadro que la
Conjetura predice.
\end{remark}

\paragraph{El \GAP\ residual, localizado.}
Lo que falta para \PR\ la Conjetura \ref{conj:hardy} / el \GAP\ (i):
\begin{enumerate}[label=(\alph*),leftmargin=2em]
\item \GAP\ \textbf{No--localidad.} $\kin$ no es $-d^2/dx^2$ local sino un generador de
salto (Lévy). Las mejoras de Hardy y las cotas CLR cl\'asicas son para operadores
locales o fraccionarios est\'andar; aqu\'i hace falta su an\'alogo para el s\'imbolo
$\Omega$ espec\'ifico. (La no--localidad --- $85\%$ del peso en saltos grandes --- fue
el obst\'aculo hist\'orico; aqu\'i reaparece, pero \emph{aislado} a esta desigualdad.)
\item \GAP\ \textbf{Alineaci\'on del extremal.} Probar \PR\ que el extremal de Hardy
es $\fp$ (hoy \NM\ por v\'ia \#9). Candidato: el extremal es el autovector de menor
autovalor de $\kin-\vv$, y la v\'ia \#9 lo identific\'o con la direcci\'on uniforme.
\end{enumerate}

% =====================================================================
\section{Estado y siguiente paso}\label{sec:status}
% =====================================================================

\begin{enumerate}[leftmargin=1.6em]
\item \PR\ (Lema \ref{lem:schur}) Reducci\'on de Schur: $A_\lambda\succeq0\Leftrightarrow$
(i)~$\kin-\vv\succeq0$ sobre $\fp^\perp$ $+$ (ii)~desigualdad escalar. Exacta, sin
hip\'otesis.
\item \PR/\NM\ (Prop.\ \ref{prop:extremal}) El d\'eficit de Hardy es el extremal
$=\fp$; la v\'ia \#9 lo confirma (rango $1$, direcci\'on--polo).
\item \OP\ (Conjetura \ref{conj:hardy}) Hardy mejorada no--local: $\kin-\vv\succeq0$
sobre $\fp^\perp$, con resto. \emph{Incondicional} (s\'olo $K,V$, no ceros). Reducida a
\GAP\ (a)~no--localidad y (b)~alineaci\'on del extremal.
\item \textbf{Consecuencia si (i) sale:} el problema entero colapsa a la \'unica
desigualdad escalar (ii) de Schur, $b+\kappa_\lambda\ge r^*M^+r$ --- una sola
comparaci\'on num\'erica entre cantidades $\sim\lambda^{-2}$, ambas computables del lado
primo.
\end{enumerate}

\noindent\textbf{Honestidad.} El \GAP\ (i) qued\'o reducido, con \'algebra exacta
(Schur) $+$ un mecanismo cl\'asico (mejora de Hardy al proyectar el extremal), a una
desigualdad funcional \emph{incondicional} cuyo \'unico obst\'aculo es la
no--localidad del s\'imbolo arquimediano. Esto \textbf{no} es RH disfrazada: es una
desigualdad de teor\'ia espectral (conteo de estados ligados de un Schr\"odinger
no--local) sobre objetos del lado primo. La batalla por (i) es ahora la batalla por una
mejora de Hardy no--local --- terreno cl\'asico, con herramientas (CLR, Lieb--Thirring,
Brezis--V\'azquez) y un obst\'aculo concreto.
\end{document}
